\documentclass[10pt,a4paper]{amsart}
\usepackage[margin=19mm]{geometry}
\usepackage{amsmath,amssymb,mathtools}
\usepackage[scr=rsfs,cal=euler]{mathalfa}
\usepackage{xfrac}
\usepackage[unicode,hidelinks,bookmarks=false]{hyperref}
\newcommand{\ie}{i.e.\ }
\newcommand{\cf}{cf.\ }
\newcommand{\resp}{respectively\ }
\newcommand{\Subsection}{\subsection}
\providecommand{\nobreakdash}{\nobreak\hbox{-}}
\setlength{\emergencystretch}{2em}
\multlinegap0pt
\usepackage{amssymb}
\usepackage{enumerate}
\newtheorem{proposition}{Proposition}
\title{Regular pullback of generalized cluster structures}
\author{Misha Gekhtman, Michael Shapiro, Alek Vainshtein}
\date{Source: arXiv:2603.12496v1 (CC BY 4.0)}
\begin{document}
\maketitle

\noindent\emph{Source and licence.} Regular pullback of generalized cluster structures. Version arXiv:2603.12496v1 (CC BY 4.0), distributed by arXiv under the Creative Commons Attribution 4.0 International licence, http://creativecommons.org/licenses/by/4.0/, as recorded in the arXiv licence metadata for this exact version and shown on the abstract page https://arxiv.org/abs/2603.12496v1. Full licence notice: \texttt{licenses/arxiv-2603.12496-CC-BY-4.0.txt}.\par\smallskip

\noindent\textit{Editorial scope.} Counted unit: the article's proposition newmu (source0.tex lines 656--673). The article's complete original proof of it is delivered with it (source0.tex lines 675--720, one \texttt{proof} environment of 46 lines). No mathematics is written, repaired or re-derived: the statement and the proof are the article's own lines, in the article's own notation, and no retained line is substituted or re-flowed.  Retained uncounted as the source-local context the proof needs: lines 590--592; lines 621--623; lines 631--633.  Nothing was omitted from any retained span, so the package carries no omission record.  No retained line contains a citation, so the wrapper carries no bibliography and no bibliography entry.  No retained line contains a figure environment, an image-inclusion command or drawing code, so no image is delivered and the package declares no asset dependency.  Editorial formatting only, in addition: an AMS \texttt{amsart} preamble that restates the article's own declarations and macros used by the retained lines, namely the environments \texttt{proposition} on separate counters, together with the standard packages those lines need.  The retained environments print in this document as Proposition 1 in order of appearance, whereas the article prints the same items with the numbering of its own full document.  The complete native source of the article as distributed in the arXiv e-print for this exact version is bundled unchanged as \texttt{sources/196281/source0.tex}.
\begin{equation}\label{cohcond}
\Lambda_k(P_m)=\max_{0\le r\le d_m}\lambda_k(T_r), \quad k\in K, 
\end{equation}

\begin{equation}\label{mutr}
\mu_{km}(r)=-\chi_{km}(r)+r\frac{\mu_{km}(d_m)-\mu_{km}(0)}{d_m}+\mu_{km}(0).
\end{equation}

\begin{proposition}\label{samefy}
For any $\Sigma$, either $\chi_{km}^\Sigma(r)=\chi_{km}(r)$, or $\chi_{km}^\Sigma(r)=\chi'_{km}(r)$, depending on the parity of the number of mutations at $m$ in a sequence connecting $\Sigma$ with the initial seed.
\end{proposition}

\begin{proposition}\label{newmu}
{\rm (i)} If $b_{mi}>0$ then
\begin{equation}\label{newmufrom}
\begin{gathered}
\mu_{km}^i(0)=\max_{0\le r\le d_m}\left\{\chi_{km}(r)-\frac r{d_m}\left(\mu_{km}(d_m)-\mu_{km}(0)\right)
+rb_{mi}\mu_{ki}(0)\right\},\\
\mu_{km}^i(d_m)-\mu_{km}^i(0)=\mu_{km}(d_m)-\mu_{km}(0)-d_mb_{mi}\mu_{ki}(0).
\end{gathered}
\end{equation}
{\rm (ii)} If $b_{im}>0$ then
\begin{equation}\label{newmutom}
\begin{gathered}
\mu_{km}^i(d_m)=\max_{0\le r\le d_m}\left\{\chi_{km}(d_m-r)-\frac r{d_m}\left(\mu_{km}(0)-\mu_{km}(d_m)\right)
+rb_{im}\mu_{ki}(d_i)\right\},\\
\mu_{km}^i(0)-\mu_{km}^i(d_m)=\mu_{km}(0)-\mu_{km}(d_m)-d_mb_{im}\mu_{ki}(d_i).
\end{gathered}
\end{equation}
\end{proposition}

\begin{proof}
(i) Recall that $\Lambda_k(P_m)=\mu_{km}(0)+\lambda_k(T_0)=\mu_{km}(d_m)+\lambda_k(T_{d_m})$. Equating the two expressions above we get
\begin{equation}\label{before}
\mu_{km}(0)%+\chi_{km}(0)
+d_mb_{mi}\lambda_k(f_i)+\lambda_k^{\mathrm {out}}
=\mu_{km}(d_m)%+\chi_{km}(d_m)
+\lambda_k^{\mathrm {in}}
\end{equation}
with
\begin{align*}
\lambda_k^{\mathrm {out}}&=d_m\sum_{\substack{l=1\\l\ne i}}^nb_{ml}\lambda_k(f_l)+\sum_{l=n+1}^sb_{ml}\lambda_k(f_l),\\
\lambda_k^{\mathrm {in}}&=d_m\sum_{l=1}^nb_{lm}\lambda_k(f_l)+\sum_{l=n+1}^sb_{lm}\lambda_k(f_l).
\end{align*}
In a similar way, after mutation at $i$ we have
\begin{multline}\label{after}
\mu_{km}^i(0)%+\chi^i_{km}(0)
+\lambda_k^{\mathrm {out}}+d_m\sum_{l=1}^nb_{mi}d_ib_{il}\lambda_k(f_l)
+\sum_{l=n+1}^sb_{mi}d_mb_{il}\lambda_k(f_l)\\
=\mu_{km}^i(d_m)%+\chi^i_{km}(d_m)
+\lambda_k^{\mathrm {in}}+d_mb_{mi}\lambda_k(f'_i).
\end{multline}
Subtracting~\eqref{after} from~\eqref{before}, 
we arrive at
\begin{multline*}
\mu_{km}^i(d_m)-\mu_{km}^i(0)=\mu_{km}(d_m)-\mu_{km}(0)\\
+d_mb_{mi}\left(d_i\sum_{l=1}^nb_{il}\lambda_k(f_l)
+\sum_{l=n+1}^sb_{il}\lambda_k(f_l)-\lambda_k(f_i)-\lambda_k(f'_i)\right).
\end{multline*}
To get the second relation in~\eqref{newmufrom}, it remains to notice that the expression in the brackets above equals
$-\mu_{ki}(0)$.

To get the first relation in~\eqref{newmufrom}, we write down~\eqref{mutr} for $\mu_{km}^i(r)$ 
taking into account that $\chi^i_{km}(r)=\chi_{km}(r)$ by Proposition~\ref{samefy},
and use condition
\[\min_{0\le r\le d_m}\mu_{km}^i(r)=0,\] 
which is equivalent to~\eqref{cohcond}. Consequently,
\[
\mu_{km}^i(0)=\max_{0\le r\le d_m}\left\{\chi_{km}(r)-\frac r{d_m}\left(\mu_{km}^i(d_m)-\mu_{km}^i(0)\right)\right\},
\]
and the first relation in~\eqref{newmufrom} follows from the second one.

(ii) The proof is similar, with~\eqref{mutr} replaced by
\[
\mu_{km}(d_m-r)=-\chi_{km}(d_m-r)+\frac r{d_m}(\mu_{km}(0)-\mu_{km}(d_m))+\mu_{km}(d_m).
\]
\end{proof}

\end{document}
